% Generated draft: requires mathematical and source review.
\documentclass[11pt]{article}
\usepackage[margin=25mm]{geometry}
\usepackage{fontspec}
\setmainfont{Latin Modern Roman}
\usepackage{amsmath,amssymb,mathtools,unicode-math}
\setmathfont{Latin Modern Math}
\usepackage{enumitem,needspace,xurl,hyperref}
\setlength{\parindent}{0pt}
\setlength{\parskip}{5pt}
\setlength{\emergencystretch}{3em}
\newcommand{\N}{\mathbb{N}}
\newcommand{\Z}{\mathbb{Z}}
\newcommand{\Q}{\mathbb{Q}}
\newcommand{\R}{\mathbb{R}}
\newcommand{\C}{\mathbb{C}}
\newcommand{\F}{\mathbb{F}}
\newcommand{\A}{\mathbb{A}}
\newcommand{\PP}{\mathbb{P}}
\newcommand{\E}{\mathbb{E}}
\newcommand{\eps}{\varepsilon}
\newcommand{\dd}{\,\mathrm{d}}
\begin{document}
\section*{UnsolvedMath ID: 701575}
\textbf{Source problem code:} catalogue entry
\subsection{Definitions and mathematical statement}
A generalized Baumslag--Solitar group is a finitely generated group \(G\) admitting an action on a tree such that every vertex stabilizer and every edge stabilizer is infinite cyclic. For the Bass--Serre formulation below, a tree is understood to be simplicial and the action to be without edge inversions; these conventions are not explicit in the supplied statement.

A finite Bass--Serre input can be encoded by a finite connected nonempty graph \(\Gamma\), with loops and multiple edges allowed, an orientation choosing one representative \(e\) of each unoriented edge, a spanning tree \(T\), and nonzero integers \(m_e,n_e\) for each chosen edge. Assign an infinite cyclic group to every vertex and edge, and let the two edge inclusions send an edge generator to the powers \(m_e\) and \(n_e\) of the generators at its initial and terminal vertices. Writing \(o(e)\) and \(d(e)\) for those vertices, the resulting group has the presentation
\[
G(\Gamma,m,n,T)=\left\langle a_v\ (v\in V\Gamma),\ t_e\ (e\in E^+\Gamma)\ \middle|\ t_e a_{o(e)}^{m_e}t_e^{-1}=a_{d(e)}^{n_e}\ (e\in E^+\Gamma),\quad t_e=1\ (e\in E^+\Gamma\cap T)\right\rangle.
\]
Here \(E^+\Gamma\) denotes the chosen oriented edge representatives. No positivity assumption is imposed on the labels.

\textbf{Question (F. A. Dudkin; described in the source as a well-known question).} Is there an algorithm which, given any two finite Bass--Serre inputs of the above form, always terminates and decides whether their presented groups are isomorphic as abstract groups?

This gives the isomorphism question a standard finite input format. The supplied source defines the class through finitely generated tree actions and states that Bass--Serre theory supplies finite presentations; it does not explicitly specify an input encoding or require the given action to have finite quotient.
\subsection{Short English statement}
Can one always decide, from two finite Bass--Serre presentations of generalized Baumslag--Solitar groups, whether the groups are isomorphic? The source defines these as finitely generated groups acting on a tree with all vertex and edge stabilizers infinite cyclic. The precise formulation uses the usual Bass--Serre action conventions, which the supplied wording leaves implicit.
\subsection{Sources}
\begin{itemize}
\item Local UnsolvedMath registry, permanent ID \(701575\), exact source problem code \texttt{catalogue entry}, title \emph{Kourovka Notebook Problem 19.27}; supplied complete source record, upstream revision \texttt{d371ee917ac9647fbb06cd6fc529deec6f3dda82}.
\item \emph{The Kourovka Notebook (Unsolved Problems in Group Theory)}, Problem 19.27, proposed by F. A. Dudkin, 19th issue (2018). The supplied statement calls this a well-known question.
\item The supplied background also identifies the 21st issue edition (2026), arXiv:1401.0300v48, \url{https://arxiv.org/abs/1401.0300v48}. The collection description names E. I. Khukhro and V. D. Mazurov as editors. These bibliographic details are supplied attribution and have not been independently verified.
\item The registry labels the problem \texttt{open} and records an update date of 2026-10-16. This is supplied registry metadata, not an independently verified mathematical or literature assessment. No separate literature-verification date is supplied.
\end{itemize}
\end{document}
